\bookchapter{A back cover with a free USB port}\label{k:cover}

We deploy code to the robot the way Berlin United does: copy the code onto a USB stick, plug the stick into the robot, done (Chapter~\ref{k:repo}). No network and no laptop are needed, which is exactly what you want in the hectic minutes before a match. On the K1 the default back cover hides the USB port. The team Ruhrbot Devils designed a replacement rear cover with an opening for the port and shared it on MakerWorld. This chapter prints and fits it, and checks that the robot sees a stick. Milestone: one (needed for USB deployment).

\begin{needbox}
\begin{itemize}[nosep]
  \item The model \emph{Booster K1 Rear Cover} by Ruhrbot Devils from MakerWorld:\\ \url{https://makerworld.com/en/models/2512402-booster-k1-rear-cover}
  \item A 3D printer, a screwdriver that fits the cover's screws, a tray for the screws.
  \item A USB stick. Background: \rtr{Devices, udev Rules and Audio}.
\end{itemize}
\end{needbox}

\begin{center}
\IfFileExists{images/k1/rear_cover.jpg}{%
  \includegraphics[width=0.6\linewidth]{k1/rear_cover.jpg}%
}{%
  \fbox{\parbox[c][5cm][c]{0.6\linewidth}{\centering\small Image of the rear cover:\\ \texttt{images/k1/rear\_cover.jpg}}}%
}
\captionof{figure}{The replacement rear cover for the Booster K1 with an opening for the USB port. Model and image: Ruhrbot Devils, ``Booster K1 Rear Cover'', MakerWorld, \url{https://makerworld.com/en/models/2512402-booster-k1-rear-cover}.}
\end{center}

\section*{Steps}
\begin{stepbox}{Step 1: get the model}
Download the model from the MakerWorld page above. Keep the link in \texttt{docs/hardware.md} of the code repository, so the next team member finds the same version.
\end{stepbox}

\begin{stepbox}{Step 2: print it}
3D print the model.
\end{stepbox}

\begin{stepbox}{Step 3: swap the covers}
Robot off. Remove the screws of the original back cover into a tray and take the cover off. If a cable is attached to the cover, note how it is plugged before you unplug it. Fit the new cover and tighten the screws only until they hold: printed plastic strips easily. Keep the original cover in the team box.
\end{stepbox}

\begin{stepbox}{Step 4: does the robot see the stick?}
Give the stick the label \texttt{K1DEPLOY} on the laptop (for a FAT stick: \verb!sudo fatlabel /dev/sdX1 K1DEPLOY!; \verb!lsblk! shows which \texttt{sdX} it is). Switch the robot on, log in with \verb!ssh k1!, and watch the kernel events while you plug the stick into the robot:
\begin{verbatim}
udevadm monitor --udev --subsystem-match=block
\end{verbatim}
Stop it with Ctrl+C and list the block devices:
\begin{verbatim}
lsblk -o NAME,LABEL,SIZE,MOUNTPOINT
\end{verbatim}
\end{stepbox}

\begin{checkbox}
\begin{itemize}[nosep]
  \item The cover sits flush and all screws hold.
  \item \texttt{udevadm monitor} prints \texttt{add} events when the stick is plugged in, and \texttt{lsblk} shows a device with the label \texttt{K1DEPLOY}.
\end{itemize}
\end{checkbox}

\begin{troublebox}
\begin{description}[leftmargin=0.6cm,nosep]
  \item[No events] Try another stick, and compare with \verb!sudo dmesg -w!, which shows the kernel's messages when a USB device appears. If nothing appears there either, check that the cable behind the cover is connected.
  \item[The label is missing] The label belongs to the partition (\texttt{sdX1}), not to the whole stick (\texttt{sdX}). Set it on the partition.
\end{description}
\end{troublebox}

\begin{keyidea}[What to remember]
\begin{itemize}[nosep]
  \item The Ruhrbot Devils rear cover frees the USB port, so code can be deployed with a stick (Chapter~\ref{k:repo}).
  \item The deploy stick has the label \texttt{K1DEPLOY}; the robot must see it in \texttt{lsblk}.
\end{itemize}
\end{keyidea}
